\chapter{数值算法与测试规范}
\label{app:numerics}

本附录把正文中的数值要求整理成可执行规范。所有阈值都应由解析误差、算法阶数或样本方差决定；本书程序内的数值只适用于其固定基准参数。

\section{推荐目录与数据流}

\begin{center}
\begin{tabular}{ll}
\toprule
目录 & 内容 \\
\midrule
\path{code/} & 可独立运行的模拟与分析程序 \\
\path{data/} & JSON/CSV 轻量指标与图数据 \\
\path{figures/} & 由代码生成的矢量或高分辨率图 \\
\path{chapters/} & 只引用稳定产物路径的正文 \\
\path{reviews/} & 逐章科学、公式、数值和编译审核 \\
\path{output/pdf/} & 可交付 PDF，不作为唯一数据存档 \\
\bottomrule
\end{tabular}
\end{center}
推荐数据流为
\[
  \text{参数清单}\longrightarrow
  \text{初态样本}\longrightarrow
  \text{轨迹/统计量}\longrightarrow
  \text{分析指标}\longrightarrow
  \text{图与表}.
\]
下游步骤不得静默修改上游参数；必要变换应生成新文件并记录来源。

\section{FFT 分步法伪代码}

对无量纲方程
\begin{equation}
  \ii\partial_t\psi=
  \left[-\frac12\partial_x^2+V+g(|\psi|^2-\delta_{\mathcal C})\right]\psi,
\end{equation}
一时间步可写为：
\begin{lstlisting}[style=bookpython]
kinetic_half = exp(-0.25j * k**2 * dt)
psi = ifft(kinetic_half * fft(psi))
local = V + g * (abs(psi)**2 - delta_c)
psi = exp(-1j * local * dt) * psi
psi = ifft(kinetic_half * fft(psi))
\end{lstlisting}
必须用实际 FFT 库验证 Parseval 约定。若将场从实空间模式振幅 $\alpha_j$ 转为幺正动量振幅，NumPy 默认变换需要额外除以 $\sqrt M$。
这里的 \texttt{delta\_c} 表示数组
$\delta_{\mathcal C}(x_j,x_j)$；只有均匀 Fourier--DVR 才能安全地把它简化为
标量 $1/\Delta x$。

\section{最低测试矩阵}

\begin{center}
\begin{tabularx}{\textwidth}{p{0.22\textwidth}XX}
\toprule
测试 & 输入 & 通过标准 \\
\midrule
真空抽样 & 独立零均值模式 & 每模 $\ensavg{|\alpha|^2}=1/2$，正规粒子数为零 \\
相干/热估计器 & 解析高斯态 & $n$ 与 $g^{(2)}$ 在统计区间内 \\
二次传播 & 谐振子或线性场 & Wigner 分布按解析辛流演化 \\
时间步长 & $\Delta t,\Delta t/2,\Delta t/4$ & 误差比符合算法阶数 \\
守恒与反演 & 时间无关闭合方程 & 范数、能量与反向误差低于预设阈值 \\
BdG 模式 & 稳定均匀模型 & 辛范数、配对与解析色散一致 \\
淬火协方差 & 稳定解析淬火 & 正常、反常矩与辛变换一致 \\
精确小系统 & Kerr 或二聚体 & 预先固定的机器可读时间窗内误差低于显式阈值 \\
截止稳定 & 至少三档物理截止 & 目标观测量存在稳定窗或报告失败 \\
\bottomrule
\end{tabularx}
\end{center}

\section{本书可运行基准}

在项目根目录使用当前 Python 环境运行
\begin{lstlisting}[style=bookpython]
python code/run_all.py
\end{lstlisting}
Windows 本机推荐显式使用项目说明中的 \texttt{AI\_group} 环境 Python。总入口依次执行：
\begin{enumerate}
  \item 第5章谐振子精确 Wigner 旋转；
  \item 第6章 Kerr 精确解与 TWA；
  \item 第8章 Bose--Hubbard 二聚体精确解与 TWA；
  \item 第11章压缩高斯协方差；
  \item 第12章分步法守恒、反演与阶数；
  \item 第13章排序估计器与 jackknife；
  \item 第14--17章均匀 Rabi 双分量场的端到端淬火、误差条、步长与截止诊断；
  \item 第15章双分量稳定淬火协方差；
  \item 第16章集体相位扩散；
  \item 第17章随机条纹相位有限尺寸标度。
\end{enumerate}
汇总写入 \path{data/test_summary.json}。端到端场案例采用均匀 Rabi 耦合，
不是 Raman SOC；第17章另一个程序是明确定义的统计玩具模型，只测试分析
公式。两者都不代表完成了具体实验参数的超固态模拟。
总入口把“进程零退出”“本轮确实产生了新 metrics”“运行来源可核验”和
“结构化科学判据通过”分成四道门。每个子程序启动前，入口先保存墙钟起点，
以及旧 metrics 的存在性、纳秒修改时间、字节数和 SHA-256；再通过环境变量
传入唯一运行 ID。子程序必须在 JSON 根层写入统一 schema 和运行记录，其中
包含运行 ID、相对脚本路径、启动前后脚本 SHA-256、启动与 metrics 定稿时间。

子进程结束后，入口要求 metrics 存在且修改时间不早于本轮墙钟起点；若旧文件
原已存在，还必须同时满足修改时间严格推进和 SHA-256 改变。因此，返回码为零
却没有写新文件的脚本不能复用旧结果过关。入口随后核对 schema、运行 ID、脚本
路径与哈希，并递归读取 \texttt{validation.checks}：既兼容以布尔值表示的简单
检查，也读取富结构及其嵌套结构中的 \texttt{passed}。由这些叶级结果重新计算的
总布尔值必须与 \texttt{validation.all\_passed} 完全一致，且该值为真，才称为
结构化基准通过。

当前清单共10项。缺失或陈旧 metrics、schema 或运行记录不匹配、没有机器可读
检查、递归重算不一致以及任一显式阈值失败，都会使总入口失败；不能用脚本零
退出替代科学判据。汇总为每项保存脚本和 metrics 的 SHA-256、metrics schema、
新旧文件状态、递归判据与时间戳，并记录 Python 与关键库版本、操作系统、Git
提交和工作区状态。

\section{收敛表模板}

不要把时间步长、模式截止和轨迹数混在同一条比较链中。最低限度应保存
三个独立的一因素扫描：
\begin{center}
\begin{tabularx}{\textwidth}{>{\raggedright\arraybackslash}p{0.13\textwidth}>{\raggedright\arraybackslash}p{0.24\textwidth}>{\raggedright\arraybackslash}p{0.21\textwidth}>{\raggedright\arraybackslash}X}
\toprule
扫描 & 固定量 & 改变量 & 判据 \\
\midrule
时间 T0--T2 & $L,M,N_s$ 与同一初态样本 &
$\Delta t,\Delta t/2,\Delta t/4$ & 差分比进入算法阶数的渐近区 \\
截止 C0--C2 & $L$、物理初态与匹配协议 &
至少三档 $M$ 或 $\Lambda$；各档时间已收敛 & 目标量在物理截止窗内稳定，或明确报告不稳定 \\
抽样 S0--S2 & 已收敛的有限模式模型与 $\Delta t$ &
$N_s,2N_s,4N_s$ & SE 近似按 $N_s^{-1/2}$ 缩小且独立种子相容 \\
\bottomrule
\end{tabularx}
\end{center}
截止扫描改变 $M$ 时也改变半量子总数，不能与单纯空间离散误差混称为
“网格收敛”。若要测试离散算符，应另做无噪声、同一带限光滑初态的空间
扫描；若要测试完整 TWA，则各截止必须重新抽样，并使用各自正确的排序
估计器。共同随机数只能在存在明确模式嵌套映射时使用，且误差应由逐轨迹
差值计算。

\section{随机数与并行}

为每个生产批次保存主种子 $s_0$，再由确定性子流函数生成轨迹种子。轨迹结果应由唯一整数 ID 标识，使失败任务可单独重跑而不改变其余随机数。并行归约时保存计数、和、平方和及必要的交叉积；浮点求和顺序变化会产生末位差异，对极端动态范围可采用成对求和或补偿求和。

共同随机数比较应保存成对轨迹 ID。差值 $D_s=X_s-Y_s$ 的误差由 $D_s$ 的样本方差计算，而不是把两个均值标准误平方相加。

\section{数据完整性与图形}

每张图至少对应一个含以下信息的数据文件：
\begin{enumerate}
  \item 横纵坐标数值和单位；
  \item 均值、误差条类型与有效样本数；
  \item 生成数据的运行 ID 与分析版本；
  \item 任何平滑、重采样、寻峰或归一化步骤；
  \item 理论曲线是解析结果、拟合还是独立模拟。
\end{enumerate}
绘图程序不得用肉眼编辑数据点。示意图与计算结果要在图注中明确区分。

\section{LaTeX 与 PDF 验收}

推荐使用 XeLaTeX 处理中文，并把构建视为测试：
\begin{enumerate}
  \item 清查未定义引用、重复标签、缺图和文献错误；
  \item 把 overfull box 视为需要定位的排版缺陷；
  \item 检查目录、书签、公式编号和交叉引用；
  \item 渲染抽查含长公式、表格、彩色框和图的页面；
  \item 保存成功构建日志、页数和稳定 PDF 哈希。
\end{enumerate}
MiKTeX 的日志目录权限警告若不影响退出码与 PDF 生成，可作为环境警告记录，但不能掩盖真正的 LaTeX 错误。
